% Standalone XeLaTeX source. Generated by assemble.py.


% ===== preamble.tex =====

\documentclass[UTF8,fontset=fandol,a4paper,openany,oneside]{ctexbook}
\usepackage{fontspec}
\setmonofont[Scale=0.80]{DejaVu Sans Mono}
\usepackage[margin=21mm,top=23mm,bottom=23mm]{geometry}
\usepackage{amsmath,amssymb,amsthm,mathtools}
\usepackage{booktabs,longtable,array}
\usepackage{xcolor,tikz}
\usetikzlibrary{arrows.meta,positioning}
\definecolor{InkBlue}{HTML}{19364D}
\definecolor{Teal}{HTML}{147D80}
\usepackage{xurl}
\usepackage[unicode,hidelinks]{hyperref}
\hypersetup{pdftitle={黎曼猜想：人类可读的证明主线},pdfauthor={基于指定Lean工程的中文数学重构},pdfsubject={ζ专用算术窗口、完整解析推导与深定理引用边界}}
\newtheorem{theorem}{定理}[chapter]
\newtheorem{lemma}[theorem]{引理}
\newtheorem{proposition}[theorem]{命题}
\newtheorem{inputtheorem}[theorem]{引用定理}
\theoremstyle{definition}
\newtheorem{definition}[theorem]{定义}
\theoremstyle{remark}
\newtheorem*{remark}{说明}
\renewcommand{\proofname}{证明}
\newcommand{\Rea}{\operatorname{Re}}
\newcommand{\Ima}{\operatorname{Im}}
\newcommand{\dd}{\,\mathrm d}
\newcommand{\Expl}[1]{\par\smallskip\noindent\textbf{这一步的意思。}#1\par\smallskip}
\setcounter{tocdepth}{1}
\setlength{\emergencystretch}{3em}
\setlength{\tabcolsep}{4pt}
\renewcommand{\arraystretch}{1.18}
\allowdisplaybreaks[1]
\raggedbottom
\urlstyle{tt}
\begin{document}
\begin{titlepage}
\vspace*{17mm}
\begin{center}
{\Large 从形式化证明到可以逐步检查的数学论证\par}
\vspace{16mm}
{\Huge\bfseries 黎曼猜想\par}
\vspace{8mm}
{\LARGE 人类可读的证明主线\par}
\vspace{12mm}
{\large 专门化到黎曼 $\zeta$ 函数\par}
{\large 算术窗口、极点与高阶导数、函数方程\par}
\end{center}
\vspace{16mm}
\noindent\rule{\textwidth}{0.5pt}
\par\medskip
本稿依据指定公开 Lean 工程及先前中文重构稿改写。正文完整展开从算术窗口估计到黎曼猜想的解析推导，并把决定性的算术几何输入准确列为引用定理；该输入的巨大上游证明没有在本稿中全部重写。
\par\medskip
\noindent\rule{\textwidth}{0.5pt}
\vfill
\noindent 整理日期：2026 年 11 月 22 日（北京时间）\par
\noindent 阅读基础：微积分、级数与复变函数的基本概念\par
\noindent 版本性质：第二版；增加完整源码汇编及算术接口展开；非作者原稿\par
\end{titlepage}
\frontmatter
\chapter*{怎样阅读这份证明}
\addcontentsline{toc}{chapter}{怎样阅读这份证明}
这份稿件把问题缩减为一个可以记住的矛盾：\textbf{如果临界线右边真的有零点，那么一个经过变换的函数就有极点；这个极点迫使某些高阶导数变大，而算术上构造出的好窗口却不允许它们那么大。}

正文不要求读者先理解整个一般 Dirichlet $L$ 函数理论。只取模 $1$ 的主特征，最终要研究的函数就是 $\zeta$。这会删去目标特征族的许多记号，但不会删去产生窗口的上游算术几何工作。

\textbf{本稿的证明范围必须准确理解：}第\ref{chap:arithmetic}章引用原工程的算术窗口定理；第\ref{chap:window}章完整证明该窗口怎样排除零点；第\ref{chap:finish}章完成参数选择与函数方程的推导。因此，本稿是一篇引用明确深定理的数学证明说明，\emph{不是}把全部上游引理都从基础重新证明的自足专著。把一项关键输入写成“引理”不会自动完成它的证明；它的来源、量词与本稿尚未展开的部分均在正文中列明。

截至本次查阅，Lean 与 Mathlib 社区于 2026 年 11 月 21 日发布的联合声明，已确认提交的命题与形式化证明，并报告命题比较与多个内核检查。该声明晚于前一版视频所采用的 MIT 12:55 UTC 状态快照。此项外部核验与本稿的人类文字重构是两件分别记录的工作；本次没有重新运行全工程核验。源码汇编恢复的是原项目各模块的文本；模块哈希一致性核验与 Lean 内核核验分别记录。原作者的四页公告也说明，论证很长，仍需进一步整理、理解和简化。参见文末资料[1]--[3]。

阅读时，可以先读每章的解释段和最后一章，掌握全链条；再回到第\ref{chap:window}章检查真正承担矛盾的每个不等式。全文不把有限个数值零点当作证明，也不把单侧上界偷换成绝对值上界。
\tableofcontents
\mainmatter


% ===== overview.tex =====

\chapter{究竟要证明什么}
\label{chap:overview}
\section{把命题写准确}
当复数 $s$ 的实部大于 $1$ 时，黎曼 $\zeta$ 函数定义为
\[
\zeta(s)=\sum_{n=1}^{\infty}n^{-s}
=\prod_{p\ \text{为素数}}(1-p^{-s})^{-1}.
\]
这些级数与乘积在这里绝对收敛。随后把 $\zeta$ 亚纯延拓到复平面；$s=1$ 是单极点，不能作为普通有限函数值参与零点讨论。

若 $\zeta(\rho)=0$，则称 $\rho$ 为零点。负偶数
$-2,-4,-6,\ldots$ 是平凡零点。本稿的目标为
\begin{theorem}[黎曼猜想的精确陈述]
\label{thm:RH}
若 $\rho\in\mathbb C$，$\rho\ne1$，$\zeta(\rho)=0$，且
$\rho\notin\{-2,-4,-6,\ldots\}$，则
\[
\Rea\rho=\frac12.
\]
\end{theorem}
这个命题说的是所有符合条件的零点。计算出很多零点在临界线上，不能替代全称证明。

\section{证明只需要两条主线相遇}
第一条主线来自算术。它构造一个记录素数与素数幂的实函数 $G(x)$，并在任意远的某些完整区间上，证明
\[
G(x)\le x^{1/2+K\varepsilon}.
\]
这里 $K$ 先固定，$\varepsilon>0$ 后选择，可以任意小。

第二条主线来自复分析。它证明，只要任意远处都有完整区间满足
$G(x)\le x^\theta$，就不可能存在 $\Rea\rho>\theta$ 的 $\zeta$ 零点。

把两条主线接起来，就能把排除区域的边界推到 $1/2$。最后利用函数方程，临界线左侧的非平凡零点会对应到已经排除的右侧，因而也不可能存在。

\[
\boxed{\text{算术窗口}}
\ \Longrightarrow\ 
\boxed{\Rea\rho\le\tfrac12}
\ \xRightarrow{\text{函数方程}}\ 
\boxed{\Rea\rho=\tfrac12}.
\]

\section{本稿保留哪些标准基础}
下文使用 $\zeta$ 的解析延拓、在 $1$ 的 Laurent 展开、Euler 乘积和函数方程；还使用 Gamma 函数的零点与极点性质、Cauchy 积分公式、恒等定理，以及全纯函数的零点孤立性。这些是已建立的复分析与特殊函数基础，可分别在所附 Mathlib 源码和资料[4]--[5]中核对。

决定性的\emph{新增算术输入}是第\ref{chap:arithmetic}章的窗口定理。它不是从上述标准基础中直接推出的初等估计，本稿不会把它隐藏在“容易证明”之类的句子中。

\section{先用一个简单函数理解极点和高阶导数}
复数 $s=\beta+it$ 对应平面上的点 $(\beta,t)$；$\Rea s=\beta$ 是横坐标。临界线 $\Rea s=1/2$ 因而是一条竖直线。说零点在临界线上，就是说它们的横坐标全部等于 $1/2$，并没有说虚部也相同。

在一个区域内可作复微分的函数称为全纯函数。若某点附近可以写成
\[
f(s)=\frac{c}{s-\rho}+h(s),\qquad c\ne0,
\]
其中 $h$ 在该点附近全纯，则称 $\rho$ 为 $f$ 的单极点。$c/(s-\rho)$ 是趋近该点时不能被有限光滑部分消去的发散项。允许有孤立极点的函数称为亚纯函数。

先只看最简单的模型 $f(s)=c/(s-\rho)$。连续求导给出
\[
f^{(k)}(z)=\frac{(-1)^k k!c}{(z-\rho)^{k+1}},
\qquad
\left|\frac{(-1)^k}{k!}f^{(k)}(z)\right|
=\frac{|c|}{|z-\rho|^{k+1}}.
\]
因此，去掉阶乘以后，第 $k$ 阶导数的大小仍记住了观察点到极点的距离。若极点距离为 $d$，而其他部分在更大半径 $r>d$ 的圆盘内全纯，那么两部分在第 $k$ 阶的尺度分别为 $d^{-k}$ 和 $r^{-k}$。它们的比值中出现
\[
\left(\frac dr\right)^k\longrightarrow0.
\]
这就是第\ref{chap:window}章中极点项最终不能被其余部分抵消的原因。正式证明还要把多个极点分离，并同时控制实轴上的导数，不能只用这个模型替代全部论证。

\section{阅读时可以反复查阅的记号表}
\begin{longtable}{@{}p{.20\textwidth}p{.73\textwidth}@{}}
\toprule
记号 & 在证明中的作用\\
\midrule
\endhead
$\rho=\beta+it$ & 假想的 $\zeta$ 零点；$\beta$ 是需要限制的实部。\\
$G(x)$ & 由素数幂加权和得到的归一化实函数，负责承接算术估计。\\
$K,\varepsilon$ & $K$ 先固定；再取小的正数 $\varepsilon$，让指数接近 $1/2$。\\
$Y,[Y,Y^{3/2}]$ & 好窗口的起点和完整区间；窗口可以出现在任意远处。\\
$U=\log Y$ & 对数坐标下的好尺度；窗口变为 $[U,3U/2]$。\\
$\Phi(s)$ & $G$ 的 Mellin 变换，也是 $g(u)=G(e^u)$ 的 Laplace 变换。\\
$\theta$ & 暂时要排除的零点边界，最后选择在 $1/2$ 与假想零点之间。\\
$\sigma,k$ & 观察点的实部及导数阶数；本稿取 $336$ 与 $\lceil420U\rceil$。\\
$d,r,q$ & 极点距离、全纯估计半径及比值 $q=d/r<1$。\\
\bottomrule
\end{longtable}


% ===== bridge.tex =====

\chapter{怎样把素数信息变成零点信息}
\label{chap:bridge}
\section{先把素数计数变得平滑}
定义冯\,芒戈尔特函数
\[
\Lambda(n)=
\begin{cases}
\log p,&n=p^a,\ p\text{ 为素数},\ a\ge1,\\
0,&\text{其他情况}.
\end{cases}
\]
例如 $\Lambda(2)=\log2$、$\Lambda(4)=\log2$、$\Lambda(8)=\log2$，而 $\Lambda(6)=0$。它把素数幂保留下来，正好与 Euler 乘积的对数导数配合。

对实数 $x\ge1$，令
\begin{align}
S(x)&=\sum_{1\le n\le x}\Lambda(n)\left(\frac{x}{n}-1\right),\label{eq:S}\\
G(x)&=S(x)-x\log x+(1+\gamma)x,\label{eq:G}
\end{align}
其中 $\gamma$ 是 Euler 常数，$\log$ 为自然对数。求和范围中的 $n$ 是整数。新项在 $x=n$ 时的权重为零，所以跨过一个素数幂时不会突然给 $S$ 加一个非零跳跃。更精确地，在任意紧的 $x$ 区间内，$S$ 是有限多个连续函数 $\Lambda(n)\max\{x/n-1,0\}$ 的和；因此 $S$、$G$ 及 $g(u)=G(e^u)$ 都连续。

例如
\[
S(4)=\log2+\frac13\log3;
\]
虽然 $4$ 是素数幂，它在此处的权重 $4/4-1$ 恰好为零。

\Expl{平滑不只是为了把图画得好看。权重 $x/n-1$ 的积分可以精确计算，随后会产生一个分母 $s(s-1)$。减去的主项则消掉本来由 $\zeta$ 在 $1$ 的极点造成的干扰。}

\begin{lemma}[无需 RH 的粗增长界]
\label{lem:growth}
存在常数 $C_0>0$，对所有 $x\ge1$ 有 $|G(x)|\le C_0x^2$。
\end{lemma}
\begin{proof}
因为 $0\le\Lambda(n)\le\log n\le n$，求和项非负，且
\[
0\le S(x)\le x\sum_{1\le n\le x}\frac{\Lambda(n)}n
\le x\lfloor x\rfloor\le x^2.
\]
再由 $\log x\le x$ 和 $x\le x^2$，
\[
|G(x)|\le x^2+x^2+(1+|\gamma|)x^2
=(3+|\gamma|)x^2.
\]
可以取 $C_0=3+|\gamma|$。
\end{proof}

\section{亲手算出 Mellin 变换}
先在 $\Rea s>2$ 定义
\[
\Phi(s)=\int_1^\infty G(x)x^{-s-1}\dd x.
\]
上面的粗界保证绝对收敛。对每个正整数 $n$，
\begin{align*}
\int_n^\infty\left(\frac{x}{n}-1\right)x^{-s-1}\dd x
&=\frac1n\frac{n^{1-s}}{s-1}-\frac{n^{-s}}s\\
&=\frac{n^{-s}}{s(s-1)}.
\end{align*}
Euler 乘积在绝对收敛区域给出
\[
-\frac{\zeta'(s)}{\zeta(s)}
=\sum_{n=1}^\infty\frac{\Lambda(n)}{n^s}.
\]
交换求和与积分的绝对收敛性可由
$\sum\Lambda(n)n^{-\Rea s}<\infty$ 及刚算出的实参数积分检验。
同时，
\[
\int_1^\infty x^{-s}\dd x=\frac1{s-1},\qquad
\int_1^\infty(\log x)x^{-s}\dd x=\frac1{(s-1)^2}.
\]
代入定义得到
\begin{equation}
\boxed{\displaystyle
\Phi(s)=-\frac{\zeta'(s)}{s(s-1)\zeta(s)}
-\frac1{(s-1)^2}+\frac{1+\gamma}{s-1}.}
\label{eq:Phi}
\end{equation}
这个等式先在收敛区域成立，右边随后给出 $\Phi$ 的亚纯延拓。

\section{为什么 \texorpdfstring{$s=1$}{s=1} 不留下一个假极点}
$\zeta$ 在 $1$ 的展开为
\[
\zeta(s)=\frac1{s-1}+\gamma+O(s-1).
\]
令 $h(s)=(s-1)\zeta(s)$，在 $1$ 延拓为 $h(1)=1$。则
$h'(1)=\gamma$，且 $h$ 在 $1$ 的一个邻域内没有零点。式\eqref{eq:Phi}可改写为
\begin{equation}
\Phi(s)=-\frac{h'(s)/h(s)+1-(1+\gamma)s}{s(s-1)}.
\label{eq:Phi-regularized}
\end{equation}
分子在 $s=1$ 等于 $\gamma+1-(1+\gamma)=0$，所以分子含有因子 $s-1$。除去它后，$\Phi$ 在 $1$ 是全纯的。

\Expl{如果不先消掉这一项，后面的“发现极点就发现零点”会被 $\zeta$ 本来的极点干扰。式\eqref{eq:G}中的校正项正是为此选择的。}

\section{每个零点都留下一个真正的极点}
设 $\rho\ne0,1$ 是 $\zeta$ 的 $m$ 重零点，$m\ge1$。在它的邻域写
\[
\zeta(s)=(s-\rho)^m a(s),\qquad a(\rho)\ne0.
\]
于是
\[
\frac{\zeta'(s)}{\zeta(s)}=
\frac m{s-\rho}+\frac{a'(s)}{a(s)}.
\]
代入式\eqref{eq:Phi}，得到
\begin{equation}
\Phi(s)=\frac{c_\rho}{s-\rho}+\text{一个在 $\rho$ 附近全纯的函数},
\qquad c_\rho=-\frac{m}{\rho(\rho-1)}\ne0.
\label{eq:residue}
\end{equation}
因此，在开右半平面 $\Rea s>0$，$\Phi$ 的极点恰好就是 $\zeta$ 的零点，且每个都是单极点。零点即使有较高重数，经过对数导数后也成为单极点，重数只进入非零系数 $c_\rho$。

\section{为何可以省掉实极点的特殊分支}
\begin{lemma}
\label{lem:no-real-poles}
$\Phi$ 在正实轴上没有极点。
\end{lemma}
\begin{proof}
当实数 $r>1$ 时，$\zeta(r)=\sum n^{-r}>0$。当 $0<r<1$ 时，考虑配对后的级数
\[
\eta(r)=\sum_{n=1}^\infty\bigl((2n-1)^{-r}-(2n)^{-r}\bigr)>0.
\]
每一项严格为正，且级数收敛。更一般地，在任何紧集 $\Rea s\ge a>0$、$|s|\le M$ 上，
\[
\left|(2n-1)^{-s}-(2n)^{-s}\right|
=\left|\int_{2n-1}^{2n}s t^{-s-1}\dd t\right|
\le M(2n-1)^{-a-1}.
\]
所以配对级数局部一致收敛，定义 $\Rea s>0$ 中的全纯函数。在 $\Rea s>1$ 将偶数项分离，可得
\[
\eta(s)=(1-2^{1-s})\zeta(s).
\]
由恒等定理，该关系延续到 $\Rea s>0$ 且 $s\ne1$。对 $0<r<1$，分母 $1-2^{1-r}<0$，故
\[
\zeta(r)=\frac{\eta(r)}{1-2^{1-r}}<0.
\]
这排除了所有正实零点；$s=1$ 在 $\Phi$ 中又已可去，因此没有正实极点。
\end{proof}

最后令
\begin{equation}
g(u)=G(e^u),\qquad u\ge0.
\label{eq:g}
\end{equation}
换元 $x=e^u$ 后，
\begin{equation}
\Phi(s)=\int_0^\infty g(u)e^{-su}\dd u,
\qquad |g(u)|\le C_0e^{2u},\qquad \Rea s>2.
\label{eq:Laplace}
\end{equation}
这样，零点已被转化为一个有实值输入的拉普拉斯变换的极点。


% ===== arithmetic.tex =====

\chapter{原 Lean 工程提供了怎样的算术信息}
\label{chap:arithmetic}
\section{把最困难的输入明确写出来}
\begin{inputtheorem}[算术几何好窗口，$\zeta$ 专门化]
\label{thm:arithmetic}
对于式\eqref{eq:G}中的 $G$，存在常数 $K>0$ 与 $\varepsilon_0>0$，使
\begin{equation}
\begin{split}
&\forall\varepsilon\in(0,\varepsilon_0)\quad\forall T\ge1
\quad\exists Y\ge T\quad\forall x\in[Y,Y^{3/2}],\\[-2pt]
&\hspace{38mm}G(x)\le x^{1/2+K\varepsilon}.
\end{split}
\label{eq:arith-window}
\end{equation}
\end{inputtheorem}

\textbf{引用范围。}这个定理由原工程的自守、Galois 上同调、宿主构造、几何链及显式公式估计组合而来。此处引用其形式化证明，并在本章展开模长度到求和界、有限特征正交、误差相加以及全部量词与阈值的拼接；这些上游构造的全部证明没有在本稿内展开。后文对\emph{这个明确输入如何推出 RH}给出完整解析论证。

原工程实际先给出固定的实数 $e_D,e_L$ 和 $\varepsilon_1>0$，指数为
$1/2+(e_D+e_L+7)\varepsilon$。取
\[
K=|e_D+e_L+7|+1,\qquad \varepsilon_0=\varepsilon_1,
\]
就得到上面的稍弱版本。由于 $x\ge1$ 且 $\varepsilon>0$，增大指数不会使上界失效。依赖于 $\varepsilon$ 的阈值可通过把 $Y$ 的起点增大而吸收。

\section{这句话中的量词各有什么作用}
\begin{enumerate}
\item $K$ 在 $\varepsilon$ 之前确定。因此 $K\varepsilon$ 可以真正趋于零。若允许任意 $K(\varepsilon)$，例如 $K(\varepsilon)=1/\varepsilon$，这个结论就没有了。
\item $Y$ 可以随着 $\varepsilon$ 和起点 $T$ 改变。不需要找一个对所有 $\varepsilon$ 都适用的共同窗口。
\item 控制的是从 $Y$ 到 $Y^{3/2}$ 的每一个实数 $x$，不是几个零散好点。
\item 结论仅为 $G(x)\le x^{1/2+K\varepsilon}$。它没有给出 $|G(x)|$ 的同样界，也不要求 $G(x)\ge0$。
\end{enumerate}

在对数坐标 $u=\log x$ 中，整个区间成为
\[
[U,3U/2],\qquad U=\log Y.
\]
它的长度是 $U/2$，随窗口向远处移动而增长。这给后面的积分核留下了足够宽的空间。

\begin{center}
\begin{tikzpicture}[x=1.6cm,y=1cm,>=Latex]
\fill[Teal!14] (2,0) rectangle (5,1.05);
\draw[->] (0,0)--(7,0) node[right] {$u=\log x$};
\draw[Teal,thick] (2,0)--(2,1.05) (5,0)--(5,1.05);
\node[below] at (2,0) {$U$};
\node[below] at (5,0) {$3U/2$};
\node at (3.5,.77) {整段成立};
\node at (3.5,.35) {$g(u)\le e^{\theta u}$};
\draw[<->] (2,1.4)--(5,1.4) node[midway,above] {长度 $U/2$};
\end{tikzpicture}
\end{center}

\section{为什么只有 $\zeta$，仍然需要深层几何}
一般工程处理有限 Dirichlet 特征族。专门取模 $1$ 的主特征时，它生成的循环群只有一个元素，群阶 $o=1$，最终保留族只剩主特征。因此目标函数确实缩减为式\eqref{eq:G}的 $G$。

但是，上游宿主中的辅助特征仍然存在。构造还使用辅助素数、分圆域、几何链与 Galois 上同调长度。特别是辅助特征的误差项不能因为最终目标只剩 $\zeta$ 就被设为零。


\section{上同调长度窗口究竟控制什么}
\label{sec:arith-length}
把数域记为 $E$，素数记为 $\ell$。对 $\ell$ 上方的一组素位及相容局部数据，源码考虑一个全局一阶 Galois 上同调群到相应局部上同调群的限制映射，取它的核，并记其 $\mathbb Z_\ell$ 模长度为 $\lambda_{\ell,j}$。所用系数是离散 Tate 扭曲 $(\mathbb Q_\ell/\mathbb Z_\ell)(-j)$。这些长度初始允许取无穷；有限性是需要证明的内容。

源码的 \texttt{VVle}$(E,t,B)$ 同时要求：
\begin{enumerate}
\item 对每个素数 $\ell$ 和 $1\le j<t$，$\lambda_{\ell,j}$ 有限；
\item 对每个有限素数集合 $S$，都有
\begin{equation}
\sum_{\ell\in S}(\log\ell)\sum_{j=1}^{t-1}\lambda_{\ell,j}\le B.
\label{eq:arith-vv-partial}
\end{equation}
\end{enumerate}
上述条件对源码允许的代数闭包和相容局部数据全称量化。它首先约束有限部分和；不能在没有说明收敛性时把定义直接写成无穷和。

\begin{lemma}[有限部分和怎样变成受控的无穷和]
设 $a_n\ge0$，且每个有限集合 $S\subset\mathbb N$ 都满足 $\sum_{n\in S}a_n\le B$。则 $\sum_{n=1}^{\infty}a_n$ 收敛，并且不超过 $B$。
\end{lemma}
\begin{proof}
取 $s_N=\sum_{n=1}^N a_n$。因为 $a_n\ge0$，$s_N$ 单调递增；有限部分和假设给出 $s_N\le B$。实数的单调有界收敛性给出极限 $s\le B$，而这正是非负项级数的和。
\end{proof}

在 \path{DivinNOfHypV.lean} 的 \texttt{vvRecord} 中，源码先核对内部记录的模长度与上述 $\lambda_{\ell,j}$ 相等，再把非负项 \texttt{vvTerm} 的有限部分和控制在 $B$ 以内。上面的引理解释了它为何同时得到
\[
\sum_{\ell}\mathrm{vvTerm}(\ell)\le B,
\qquad \sum_{\ell}\mathrm{vvTerm}(\ell)\text{ 收敛}.
\]
这里模长度的记录相等依赖对应的数据结构；本段展开的是由非负有限部分和界得到可求和性的数学推理。

\paragraph{原始窗口的大小。}
\path{ExistsPageN.lean} 中的 \texttt{VWinAt} 固定指数 $401$，对
$0<\varepsilon<1/52$ 及满足
\[
x^{2+8\varepsilon}\le2^{e(x)}\le2x^{2+8\varepsilon}\qquad(x\ge1)
\]
的整数值函数 $e(x)$，给出常数 $J\ge1$ 和 $X_0$。对足够大的 $X$ 以及满足指定同余、大小和整除条件的辅助素数对 $q,d$，存在
\[
X^{1/J}\le Y,\qquad Y^2\le X,
\]
使所有 $x\in[Y,Y^2]$ 上相应分圆特征域满足
\begin{equation}
\mathrm{VVle}\left(E,\lfloor x^{401}\rfloor,
\sqrt{\exp(x^\varepsilon)}\right).
\label{eq:arith-vwin}
\end{equation}
例如素数对条件包含 $q\equiv3\pmod8$、$d\equiv7\pmod{16}$、$q\mid d+1$，以及指定范围内的其他奇素数也整除 $d+1$。完整条件保留在源码 \texttt{WinPair} 中；这条窗口定理本身并不替代辅助素数对的存在性构造。

\Expl{这里首先得到的是某些上同调群的总长度预算。它还没有直接出现素数平滑函数 $G$。后面的宿主和几何比较负责把这份预算转成 $G$ 的上界。}

\section{数据族、宿主和统一常数}
\label{sec:arith-family}
对固定的小参数 $\varepsilon$，源码使用一个数据族 $\mathcal F$。在尺度 $x$ 上有一组索引 $I_x$；每个 $i\in I_x$ 带一个算术宿主 $H_{x,i}$ 和尺寸数据。令 $\mathcal W$ 表示该族的好尺度集合。窗口性质写为
\begin{equation}
\forall T\quad\exists Y\ge T,
\qquad [Y,Y^\lambda]\subseteq\mathcal W,
\qquad\lambda=\frac32.
\label{eq:arith-family-windows}
\end{equation}
这里的宿主是承载数域、有限群、特征及相容映射的数据；它的作用是使上同调长度、几何量和特征和可以在同一组对象上比较。

\paragraph{宿主的存在。}
\path{SupplyChi.lean} 的实际接口是
\[
\texttt{VWinAt}(\psi,401)
\ \Longrightarrow\
\texttt{HostExW}\bigl(3/2,(\texttt{GChi}\ \psi).\texttt{req}\bigr).
\]
其证明调用 \texttt{supplyOfAdm}，并把 \texttt{admissibleClassesTwistedChar} 提供的可容许类作为输入。这是一个带有明确前提和辅助构造的宿主存在结论。原始长度窗口 $[Y,Y^2]$ 与最终族窗口 $[Y,Y^{3/2}]$ 属于这条链的不同层次，不能混用。

\paragraph{展开 \texttt{HostExW} 的量词。}
该性质的定义位于 \path{ThinScalesIffHostBound.lean}。对于每个 $0<\varepsilon<1/52$，它先选取
\[
C_{T_2},C_8,C_5,c_6,C_b,x_0
\quad\text{及一个尺度系统 }\mathcal S=(J,\mathcal W),
\]
其中 $C_8\ge1$、$C_5,c_6,C_b>0$、$J\ge1$，$\mathcal W$ 具有式\eqref{eq:arith-family-windows}的窗口性质。然后，对每个与 $\varepsilon,J$ 相容、满足 $x\ge x_0$ 且 $x\in\mathcal W$ 的尺度记录，才存在秩数据 $R$、宿主 $H$ 及辅助素数 $q_0$，使
\[
\mathrm{RankVal}(s,R),\qquad
H.C_{T_2}=C_{T_2},\qquad
\mathrm{HostGood}(s,C_8,C_5,c_6,C_b,H,q_0),\qquad P(C_5,s,R,H,q_0)
\]
同时成立。这里 \texttt{HostGood} 还包括素数同余、厚薄部分估计、分离性及指定水平素数集合等要求；其全部构造不由“存在宿主”这几个字自动完成。

对当前特征 $\psi$，\path{Supply.lean} 实际把请求 $P$ 定义成
\begin{equation}
\begin{split}
&\mathrm{EAtH}\bigl(\mathrm{Level.cyc}(\psi),H,T,x,\varepsilon,4\bigr)\\
&\hspace{16mm}\land\ \mathrm{VVle}(E_H,t,\sqrt b).
\end{split}
\label{eq:arith-gchi-request}
\end{equation}
因而误差条件与长度条件是在这一步一起要求并一起取得的。\texttt{FReq.cls} 再把它们沿数据族的所有宿主保存下来。

\paragraph{为何把平方窗口缩成 $3/2$ 次方窗口。}
\path{WinPairs.lean} 中的 \texttt{winPairs} 给出一个可以直接检查的缩放步骤。令 $c\ge1$ 是分圆域构造带来的固定尺度因子，取
\[
\alpha=2+8\varepsilon>0,\qquad k=c^{1/\alpha}\ge1.
\]
将原始平方窗口的起点记为 $W$，并把它选得足够大，使
$W\ge\max\{1,k^3,k^{400}\}$。令 $Z=kW$。由 $k^3\le W$，
\[
Z^{3/2}=k^{3/2}W^{3/2}\le W^{1/2}W^{3/2}=W^2.
\]
所以对每个 $x\in[Z,Z^{3/2}]$，记 $x'=x/k$，便有
\[
W\le x'\le x\le W^2.
\]
原始长度窗口因此可以在 $x'$ 上使用。整数参数的尺度界也正确变换：
\[
(x')^\alpha\le2^{e(x')}\le2(x')^\alpha
\ \Longrightarrow\
x^\alpha\le c\,2^{e(x')}\le2x^\alpha.
\]
此外，$k^{400}\le W\le x'$ 给出
\[
x^{400}=k^{400}(x')^{400}\le(x')^{401},
\qquad
\sqrt{\exp((x')^\varepsilon)}\le\sqrt{\exp(x^\varepsilon)}.
\]
\texttt{VVle} 对长度截断参数向下、对上界参数向上具有单调性：有限部分和只会减少，允许的预算只会增大。因此可从原始 $\lfloor(x')^{401}\rfloor$ 长度界得到新尺度的 $\lfloor x^{400}\rfloor$ 长度界。原工程随后用秩估计 $R.t\le\lfloor x^{400}\rfloor$ 得到式\eqref{eq:arith-gchi-request}所需的截断长度。

这里完整解释了 $2\to3/2$ 的窗口缩放和 $401\to400$ 的参数余量。辅助素数对的选择、特征域的构造及秩估计本身，仍引用 \texttt{winPairs}、\texttt{supplyOfAdmOf} 所调用的相应定理。

\paragraph{同一个族必须同时满足两类要求。}
对于满足 $(\texttt{GChi}\ \psi).\texttt{cls}(\lambda,\mathcal F)$ 的族，\path{ZeroFreeAllOfTopG.lean} 分别提取出：
\begin{enumerate}
\item \texttt{HypV}：超过某个阈值，对所有 $i\in I_x$，宿主上的 \texttt{VVle} 长度界成立；
\item \texttt{EFamC}：超过另一个阈值，对所有 $i\in I_x$，宿主满足后文的显式公式误差条件 \texttt{EAtH}。在这一步，源码选择的误差常数为 $4$。
\end{enumerate}
两个结论来自同一个族的类成员证明，因此可以在同一个宿主上联合使用。若分别找出两个互不相干的族，不能直接把它们的结论拼在一起。

\paragraph{几何常数的先后顺序。}
\texttt{FamilyClassConnected} 的完整量词骨架为
\begin{equation}
\begin{split}
&\exists e_D,e_L\in\mathbb R\quad\exists\varepsilon_1>0\quad
\forall\varepsilon\in(0,\varepsilon_1)\quad\exists\mathcal F,\\
&\exists\kappa,c_{HH},e_\kappa,C_L,x_T\quad
\forall x\ge x_T,\quad
x\in\mathcal W\Longrightarrow
\mathrm{ChainTied}(\mathcal F,e_D,e_L,\kappa,c_{HH},e_\kappa,C_L,x).
\end{split}
\label{eq:arith-tied-quantifiers}
\end{equation}
这里省去的只是“$\mathcal F$ 属于指定类”这一并列条件，仍须保留。$e_D,e_L,\varepsilon_1$ 在选 $\varepsilon$ 之前固定；$\kappa,c_{HH},e_\kappa,C_L,x_T$ 则允许依赖 $\varepsilon$ 和选出的族。式\eqref{eq:arith-window}中的固定指数系数只取决于前一组常数。

\Expl{不能把所有常数都说成与 $\varepsilon$ 无关；源码没有这个要求。需要统一的恰好是最后乘在 $\varepsilon$ 前面的指数系数。其余常数可以变大，因为稍后允许把好窗口选得更远。}

\paragraph{为什么九类几何条件可以在同一个族上成立。}
这里有一个容易在口头叙述中遗漏的量词问题。源码的 \texttt{Closed}$(\mathcal C,P)$ 并不只是说“某个族满足 $P$”，而是
\[
\exists\delta>0\quad\forall\varepsilon\in(0,\delta]
\quad\exists\mathfrak K\quad
\forall(\mathcal F,C_H),\quad
\mathcal C(\mathcal F)\land\mathrm{Large}(\mathcal F,\mathfrak K,C_H)
\Longrightarrow P(\mathcal F,C_H).
\]
$\mathfrak K$ 是若干阈值组成的记录，源码称 \texttt{TransportedAlt}。两个这样的记录可以逐项取较大值 $\mathfrak K\vee\mathfrak K'$；满足较大记录就同时满足二者。故若有两个 \texttt{Closed} 条件，其共同小参数阈值取 $\min\{\delta,\delta'\}$，共同大小记录取 $\mathfrak K\vee\mathfrak K'$，就得到它们的合取。有限多个条件可以重复这一过程。

与之相配的 \texttt{DesignExists} 是
\[
\exists\delta_D>0\quad\forall\varepsilon\in(0,\delta_D]
\quad\forall\mathfrak K\quad\exists(\mathcal F,C_H),
\]
使指定的类成员条件、$\mathrm{Large}(\mathcal F,\mathfrak K,C_H)$、所选参数条件及两类尺寸条件同时成立。注意其中是“对每个阈值记录，都存在一个族”。

\path{DesWin.lean} 的 \texttt{topDesWin} 先把九类几何相容性条件合并为一个 \texttt{Closed} 条件，再取它与 \texttt{DesignExists} 的小参数阈值的最小值。固定这样的 $\varepsilon$ 后，\emph{先由合并条件选出 $\mathfrak K$，再把这个 $\mathfrak K$ 送入设计存在性定理，选出同一个族 $\mathcal F$}。这样九类条件确实同时作用于这一个族。最后结合长度条件与秩关系得到所需的几何链。

本段证明的是条件合并和选择顺序的有效性；九类条件的各自证明、以及设计存在性定理的构造内容，仍由原工程中的相应几何定理承担。

\paragraph{从长度预算到几何误差的一步吸收。}
\path{DivinNOfHypV.lean} 还明确证明：当 $\varepsilon>0$ 时，充分大的 $x$ 满足
\[
(\log x)^2\le x^\varepsilon.
\]
例如令 $t=\log x>0$，由 $e^{\varepsilon t}\ge(\varepsilon t)^3/6$，当 $t\ge6/\varepsilon^3$ 时就有 $e^{\varepsilon t}\ge t^2$。因此这类对数平方损失可以吸收到一个 $x^\varepsilon$ 中。源码把这一点、\texttt{HypV} 提供的可求和长度预算以及指定的秩相等条件联合使用，得到名为 \texttt{DIVinN} 的几何尺寸控制。完整几何链还需要其他结构输入，不能把这一项当成整条链的全部证明。

\section{显式公式误差条件究竟包含什么}
\label{sec:arith-eath}
把源码的有限群框架记为 $\varphi$，其中有有限阿贝尔群 $A$、子群 $A_0$ 和一个指定整数 $o$。对宿主 $H$、有限采样集 $T\subseteq A$，\texttt{EAtH}$(\Lambda,H,T,x,\varepsilon,C)$ 包含以下内容：
\begin{enumerate}
\item 宿主群框架的整数 $o$ 与目标特征层 $\Lambda$ 的 $o$ 相等；
\item 应保留的目标特征及指定例外特征按规定嵌入宿主；
\item 采样点在各陪集中的分布满足误差界：令
\[
\lambda_\gamma=
\begin{cases}
4/(3o+1),&\gamma\in A_0,\\
3/(3o+1),&\gamma\notin A_0,
\end{cases}
\]
则对每个 $\gamma\in A$，
\begin{equation}
\left|\#\{t\in T:t-\gamma\in A_0\}-\lambda_\gamma\#T\right|
\le x^{-1+\varepsilon}\#T;
\label{eq:arith-shares}
\end{equation}
\item 未保留而属于源码 \texttt{Bounded} 类的辅助特征，其归一化平滑和满足点态范数界 $x^{1/2+\varepsilon}$；
\item 相应本原特征在 $1$ 处的对数导数有界于 $C\log x$。
\end{enumerate}
这是 \path{ExplicitFormulaErrorOfHostConditions.lean} 中的结构定义。它表明：辅助误差控制包括特征筛选、陪集计数和两种解析估计，并不是仅仅规定一个总误差符号。

\section{几何链怎样接到解析比较}
\label{sec:arith-coreath}
几何链条件 \texttt{ChainTied} 首先存在量化出一组相容数据。\path{LevelGmLeOnFamilyScales.lean} 的 \texttt{tiedCore\_of\_chainTied} 从中选择一个 $i_0\in I_x$，把它整理成对应宿主上的 \texttt{CoreAtH}。

该输入包中有有限索引集 $U_0$、每个 $u\in U_0$ 上非空有限集 $W_u$、顶部索引集、自然数 $m,f$，以及把 $U_0$ 映到 $A$ 的映射 $\operatorname{tr}$。它至少覆盖 $A_0$：
\[
\forall s\in A_0\quad\exists u\in U_0,
\qquad\operatorname{tr}(u)=s.
\]
其余条件把这些几何数据与实际的算术和联系起来。按作用可读成下表。

\begin{longtable}{@{}p{.29\textwidth}p{.66\textwidth}@{}}
\toprule
源码条件 & 在推导中的用途\\
\midrule
\endhead
\texttt{HostAddCharCompatible} & 对非平凡的宿主加法特征，把辅助参数 $\ell(\psi)$ 识别为相应本原 Dirichlet 特征在 $1$ 处的对数导数。\\
\texttt{IdentU} & 把几何记录中的主项和实际的有限特征求和逐项识别。\\
\texttt{SizeData₂}、\texttt{ErrWindows}、\texttt{BNBoundsU} & 给出逐索引几何差值的尺寸与误差控制。\\
\texttt{SizesA}、\texttt{ThinScales}、\texttt{CostWindows}、\texttt{RestSize} & 控制厚薄区域、尺度移动和后续平均过程中引入的剩余项。\\
$T_2\subseteq A_0$ & 使下一节的子群特征正交可以用于所选的厚区域。\\
\bottomrule
\end{longtable}

这些条件连同 $D.x=x$、$D.\varepsilon=\varepsilon$、
$C.eDiv\le e_D$、$C.eLoss\le e_L$ 等常数比较，构成了 \texttt{CoreAtH} 的精确定义。表中尚未展开的几何条件仍按其来源引用；本稿不把它们省略成没有前提的数值不等式。

\paragraph{逐项差值的识别。}
\texttt{IdentU} 给出
\[
\operatorname{Main}(u)
=\sum_{t\in T}\sum_{j=1}^{m}
\Rea\operatorname{MAIN}\bigl(\ell,a(t),w(\operatorname{tr}(u),j)\bigr).
\]
与同一记录中的算术和逐项相减，得到
\[
\operatorname{gapA}(\operatorname{tr}(u))
=\operatorname{ASK}(u)-\operatorname{Main}(u).
\]
\path{ChTop.lean} 的 \texttt{ch\_top} 先使用逐单位几何界，再用窗口误差和基数界，把右边控制为
\begin{equation}
\operatorname{gapA}(\operatorname{tr}(u))
\le x^{C_{\rm err}\varepsilon}\mathcal B(x),
\label{eq:arith-gap}
\end{equation}
其中 $\mathcal B(x)\ge0$ 是源码明确给出的尺寸误差括号。由于 $\operatorname{tr}$ 覆盖 $A_0$，这个上界对每个需要平均的 $s\in A_0$ 都成立。原工程还给出
$C_{\rm err}\le e_D+e_L+3$，从而把式\eqref{eq:arith-gap}中的幂统一放大为 $x^{(e_D+e_L+3)\varepsilon}$。

\section{有限特征正交怎样保留目标项}
\label{sec:arith-orthogonality}
令 $E_0=(\mathbb Z/2\mathbb Z)^2$。对有限阿贝尔群 $E_0\times A$ 的复特征 $\psi$，定义
\[
\Psi(A_0)=\{\psi:\ \psi(0,a)=1\text{ 对所有 }a\in A_0\}.
\]
下面把源码所用的基本正交步骤完整证明。

\begin{lemma}[子群上的有限特征和]
若 $\chi:A_0\to\mathbb C^\times$ 为特征，则
\[
\sum_{a\in A_0}\chi(a)=
\begin{cases}
\#A_0,&\chi\equiv1,\\
0,&\chi\not\equiv1.
\end{cases}
\]
\end{lemma}
\begin{proof}
平凡特征的每一项等于 $1$，第一种情况直接成立。若 $\chi$ 非平凡，选取 $a_0$ 使 $\chi(a_0)\ne1$，并令 $S=\sum_a\chi(a)$。映射 $a\mapsto a+a_0$ 置换 $A_0$，故
\[
S=\sum_a\chi(a+a_0)=\chi(a_0)\sum_a\chi(a)=\chi(a_0)S.
\]
所以 $(1-\chi(a_0))S=0$。第一个因子非零，故 $S=0$。
\end{proof}

特别地，令 $e_1(a)=((1,0),a)$。对 $s\in A_0$ 和 $t\in A$，利用群元素恒等式
\[
-e_1(s+a-t)=e_1(t)+(0,-s)+(0,-a),
\]
以及 $-1=1$ 在 $\mathbb Z/2\mathbb Z$ 中成立，得到
\begin{equation}
\sum_{a\in A_0}\psi\bigl(-e_1(s+a-t)\bigr)
=
\begin{cases}
\#A_0\,\psi(e_1(t)),&\psi\in\Psi(A_0),\\
0,&\psi\notin\Psi(A_0).
\end{cases}
\label{eq:arith-character-projection}
\end{equation}
这对应 \path{ChCore.lean} 的 \texttt{sum\_A}。几何差值在子群上求和后，不满足子群平凡性要求的特征被准确消去。

\paragraph{保留项与辅助项的分离。}
设 $\mathcal K\subseteq\Psi(A_0)$ 为要保留的有限特征集合，$P_\psi(x)$ 为与宿主相容的归一化平滑和，$w_\psi$ 为保留权重。定义
\[
v_\psi(x)=
\begin{cases}
w_\psi\Rea P_\psi(x),&\psi\in\mathcal K,\\
-|P_\psi(x)|,&\psi\notin\mathcal K.
\end{cases}
\]
将有限集合拆成 $\mathcal K$ 与它的补集，就得到完全精确的恒等式
\begin{equation}
\sum_{\psi\in\Psi(A_0)}v_\psi(x)
=\underbrace{\sum_{\psi\in\mathcal K}w_\psi\Rea P_\psi(x)}_{G_{\mathcal K}(x)}
-\underbrace{\sum_{\psi\in\Psi(A_0)\setminus\mathcal K}|P_\psi(x)|}_{\Sigma(x)}.
\label{eq:arith-g-minus-sigma}
\end{equation}
这就是源码 \texttt{sum\_vv} 的内容。因此辅助项以一个明确的非负误差和出现，不会由于目标只取 $\zeta$ 就自动消失。

\paragraph{从投影到数值界仍需尺度控制。}
正交恒等式负责筛选特征；要从式\eqref{eq:arith-gap}得到 $G_{\mathcal K}$ 的上界，还需要处理采样点 $a(t)/b$ 与 $x$ 的差异，以及厚薄区域和余项。源码 \texttt{ch\_core} 的前提明确包含
\[
\left|\frac{a(t)}b-x\right|\le1,
\qquad m\le x^{1/2+\varepsilon/2},
\qquad xm\le k,
\]
以及上一节列出的尺寸和误差条件。它通过有限求和、平滑和的移动估计及尺寸不等式，给出
\begin{equation}
G_{\mathcal K}(x)\le
x^{1/2+(C+3)\varepsilon}+\Sigma(x)
\quad\text{当式\eqref{eq:arith-gap}的指数系数为 }C\ge0.
\label{eq:arith-chcore}
\end{equation}
本节完整展开了筛选恒等式；\texttt{ch\_core} 中全部厚薄区域估计与尺寸约束的构造仍作为原工程定理引用。因而不能说“由特征正交”四个字就已经证明了式\eqref{eq:arith-chcore}的全部内容。

\section{辅助误差的有限求和与指数合并}
\label{sec:arith-error-sum}
在原工程宿主中，子群指数 $[A:A_0]=o$。\texttt{CardPsiA.lean} 中的基数关系可以这样理解：$\Psi(A_0)$ 的特征恰好经过商映射
\[
E_0\times A\longrightarrow E_0\times(A/A_0).
\]
若特征在核 $\{0\}\times A_0$ 上平凡，就可以在每个陪集上选一个代表定义商群特征；不同代表的差在核中，所以定义与代表选择无关。反过来，商群特征拉回后一定在核上平凡；满射性保证这个拉回是单射。因此两边的特征集合一一对应。

使用有限阿贝尔群“复特征的个数等于群阶”的标准对偶事实，可得
\begin{equation}
\#\Psi(A_0)=\#\widehat{E_0\times(A/A_0)}
=\#E_0\cdot[A:A_0]=4o.
\label{eq:arith-card-psi}
\end{equation}
这项标准对偶事实也可由有限阿贝尔群的循环群分解理解：阶为 $n$ 的循环群，其生成元可映到任意一个 $n$ 次单位根，恰有 $n$ 个特征；直积的特征数相乘。这里只引用有限阿贝尔群的标准结构定理，不把它混入新的算术几何假设。
\texttt{EAtH} 的逐项界及本原特征识别，保证对每个未保留特征有
$|P_\psi(x)|\le x^{1/2+\varepsilon}$。于是
\begin{align}
\Sigma(x)
&=\sum_{\psi\in\Psi(A_0)\setminus\mathcal K}|P_\psi(x)|\notag\\
&\le\#\bigl(\Psi(A_0)\setminus\mathcal K\bigr)x^{1/2+\varepsilon}\notag\\
&\le4o\,x^{1/2+\varepsilon}.
\label{eq:arith-sigma-general}
\end{align}
这是 \path{SigALe.lean} 的 \texttt{sigA\_le} 所完成的推理。特征集合大小和宿主相容性有各自的证明，以上三行是它们成立之后的完整误差合并。

令 $E=e_D+e_L$。对实际构造的数据，核心条件给出 $E\ge0$。把式\eqref{eq:arith-chcore}中的 $C$ 取为 $E+3$，得到
\begin{align}
G_{\mathcal K}(x)&\le x^{1/2+(E+6)\varepsilon}+\Sigma(x),\label{eq:arith-main-general}\\
\Sigma(x)&\le4o\,x^{1/2+\varepsilon}.
\end{align}
由于 $x\ge1$、$E\ge0$，
\[
x^{1/2+\varepsilon}\le x^{1/2+(E+6)\varepsilon}.
\]
再要求 $x^\varepsilon\ge1+4o$，便有
\begin{align*}
G_{\mathcal K}(x)
&\le(1+4o)x^{1/2+(E+6)\varepsilon}\\
&\le x^\varepsilon x^{1/2+(E+6)\varepsilon}
=x^{1/2+(E+7)\varepsilon}.
\end{align*}

专门化到本稿的模 $1$ 主特征时，$o=1$，目标层的保留集合为 $\mathrm{Level.KEPT}=\{1\}$。通过 \path{GmAEq.lean} 的 \texttt{gmA\_eq} 所给出的本原特征与权重识别，宿主上的 $G_{\mathcal K}$ 等于目标层的平滑和，即式\eqref{eq:G}中的 $G$。这里区分宿主加法特征与目标 Dirichlet 特征，避免把两种类型中的平凡元混写。因此
\begin{align}
G(x)&\le x^{1/2+(E+6)\varepsilon}+\Sigma(x),\label{eq:arith-main}\\
\Sigma(x)&\le4x^{1/2+\varepsilon},\label{eq:arith-error}\\
G(x)&\le x^{1/2+(E+7)\varepsilon}\qquad(x^\varepsilon\ge5).
\label{eq:arith-final-numeric}
\end{align}

\Expl{最后的 $7$ 分为 $3+3+1$：几何误差指数放大增加 $3$，\texttt{ch\_core} 的尺度与平均估计增加 $3$，最后用一个 $x^\varepsilon$ 吸收固定系数。这三个损失都乘在可以后选的小参数 $\varepsilon$ 上。}

\section{把全部阈值放进同一个好窗口}
\label{sec:arith-thresholds}
到这里还必须核对：各个估计是否真的在同一个 $x$、同一个宿主上成立。

固定 $\varepsilon$ 和已经选出的族 $\mathcal F$。设
\begin{enumerate}
\item $x_E$ 是 \texttt{EFamC} 的阈值：对所有 $x\ge x_E$ 和所有 $i\in I_x$，\texttt{EAtH} 成立；
\item $x_T$ 是几何链的阈值：对所有 $x\ge x_T$ 且 $x\in\mathcal W$，\texttt{ChainTied} 成立；
\item $x_C$ 是 \texttt{core\_gm} 在这一组固定常数下给出的数值阈值。它已经包括几何大尺度条件、\texttt{ch\_core} 的阈值及 $x^\varepsilon\ge1+4o$。
\end{enumerate}
令
\[
X_* = \max\{1,x_E,x_T,x_C\}.
\]
对任意 $x\ge X_*$ 且 $x\in\mathcal W$：先由几何链选取一个 $i_0\in I_x$，取得该宿主的 \texttt{CoreAtH}；再由“对所有 $i$”成立的误差条件，取得同一 $i_0$ 的 \texttt{EAtH}；最后使用 \texttt{core\_gm}，得到
\[
G(x)\le x^{1/2+(E+7)\varepsilon}.
\]
这正是 \texttt{windowInequality\_of\_artinInputs} 的三阈值最大值论证。无需让每个宿主都满足几何链；每个好尺度存在一个满足它的宿主已经足够，因为误差条件对全部宿主成立。

现给任意起点 $T\ge1$，使用式\eqref{eq:arith-family-windows}，把起点增大到 $\max\{T,X_*\}$，选取 $Y$ 使
\[
Y\ge\max\{T,X_*\},
\qquad[Y,Y^{3/2}]\subseteq\mathcal W.
\]
对这个区间内的每个 $x$，都有 $x\ge X_*$，所以整个区间都满足同一个上界。最后取
\[
K=|E+7|+1,\qquad\varepsilon_0=\varepsilon_1,
\]
就得到引用定理\ref{thm:arithmetic}。其中 $K$ 在选取 $\varepsilon$ 之前确定；$X_*$ 和 $Y$ 可以在其后确定。这一先后顺序完整保留了后文排除零点所需要的量词。

\section{本章已展开到哪一层}
本章已经把以下推理写成可逐步检查的数学过程：非负部分和界到可求和性；平方窗口缩放成 $3/2$ 次方窗口；长度截断 $401$ 到 $400$ 的余量；共同几何条件的合取与设计族的选择顺序；统一指数常数与可变阈值的区别；有限群特征正交；保留项与辅助误差的准确拆分；$4o$ 项误差求和；$E+7$ 指数合并；以及在同一宿主和同一窗口上汇合全部条件。

仍作为深定理引用的部分是：产生 \texttt{VWin} 的自守与 Galois 上同调论证；可容许宿主及相关辅助素数的构造；产生相容几何链的全部步骤；逐单位几何比较；以及 \texttt{ch\_core} 中尚未在本章完整转写的尺度与厚薄区域估计。这些事项有原工程的明确入口，不能因本章已经完成接口之间的拼接，就标为已经全部完成了人类文字证明。


% ===== window_proof.tex =====

\chapter{窗口型无零点定理}
\label{chap:window-proof}
\label{chap:window}

本章证明一个明确的解析推断：如果归一化的平滑素数和在任意远处的整段窗口上受到控制，那么相应区域内就不能有 $\zeta$ 零点。窗口估计是这一推断的输入；下面完整证明从该输入到无零点结论的过程。

沿用前文的实值函数 $G$、$g(u)=G(e^u)$ 及其拉普拉斯变换 $\Phi$。本章使用前文已经建立的以下性质：$g$ 连续，存在固定 $C_0>0$ 使
\begin{equation}
 |g(u)|\le C_0e^{2u}\qquad(u\ge0);
 \label{eq:win-growth}
\end{equation}
当 $\operatorname{Re}s>2$ 时，
\begin{equation}
 \Phi(s)=\int_0^\infty g(u)e^{-su}\,du.
 \label{eq:win-laplace}
\end{equation}
此外，$\Phi$ 亚纯延拓到半平面
\[
 \mathcal H=\{s\in\mathbb C:\operatorname{Re}s>0\}.
\]
它的极点集合 $Z$ 恰为 $\zeta$ 在 $\mathcal H$ 中的零点集合。每个 $\rho\in Z$ 都是 $\Phi$ 的单极点，留数 $c_\rho\ne0$；$\Phi$ 在正实轴上没有极点。这里 $\zeta$ 的 $s=1$ 极点已经在构造 $\Phi$ 时被消去。

\begin{theorem}[窗口型无零点定理]
\label{thm:window-zeta}
固定 $0<\theta<1$。假定
\begin{equation}
 \forall Y_0\ge1\quad
 \exists Y\ge Y_0\quad
 \forall x\in[Y,Y^{3/2}],
 \qquad G(x)\le x^\theta.
 \label{eq:win-pointwise}
\end{equation}
则 $\Phi$ 在 $\operatorname{Re}s>\theta$ 内没有极点。因此，
\[
 \operatorname{Re}s>\theta,\qquad s\ne1
 \quad\Longrightarrow\quad \zeta(s)\ne0.
\]
\end{theorem}

证明将比较同一组高阶导数的两个界。窗口估计给出上界；若存在一个待排除的极点，它又会迫使这些导数具有更大的下界。关键是把两种估计放在同一尺度上。

\section{把窗口上界变成导数上界}

令 $U=\log Y$。幂区间 $[Y,Y^{3/2}]$ 变成对数区间
\[
 W_U=[U,3U/2].
\]
式 \eqref{eq:win-pointwise} 意味着：对任意 $U_0\ge1$，总能选到 $U\ge U_0$，使
\begin{equation}
 g(u)\le e^{\theta u}\qquad(u\in W_U).
 \label{eq:win-log-window}
\end{equation}
这是因为可先对起点 $Y_0=e^{U_0}$ 使用窗口假设，再取对数。以后称满足 \eqref{eq:win-log-window} 的 $U$ 为好尺度。所需的只是好尺度无界。

对整数 $k\ge0$ 和 $\operatorname{Re}z>2$，定义
\begin{equation}
 A_k(z)=\frac{(-1)^k}{k!}\Phi^{(k)}(z)
       =\frac1{k!}\int_0^\infty g(u)u^ke^{-zu}\,du.
 \label{eq:win-Ak}
\end{equation}
第二个等式可以逐次在积分号下求导。确实，在任意这样的 $z$ 的一个小邻域内，实部都有统一下界 $\sigma_0>2$，而第 $k$ 阶被积函数的绝对值由
\[
 C_0u^ke^{-(\sigma_0-2)u}
\]
控制。它在半轴上可积，因而每一阶求导都有所需的可积控制函数。这也说明 $\Phi$ 在 $\operatorname{Re}z>2$ 全纯。

固定
\begin{equation}
 \sigma=336.
 \label{eq:win-sigma}
\end{equation}
所有好尺度都使用这个相同的实部。$336$ 只是让后面的两侧尾估计同时成立的一种方便选择；下面会直接验证这些数值，不把它们当作额外假设。对 $u\ge0$，记
\[
 w_k(u)=\frac{u^ke^{-\sigma u}}{k!}\ge0.
\]
由于 $g$ 实值，$A_k(\sigma)$ 也是实数。三角不等式给出
\begin{align}
 |A_k(\sigma+it)|+\operatorname{Re}A_k(\sigma)
 &\le \int_0^\infty \bigl(|g(u)|+g(u)\bigr)w_k(u)\,du.
 \label{eq:win-positive-part}
\end{align}
这一步是使用单侧上界的关键。对于实数 $v$，
\[
 |v|+v=2\max\{v,0\}.
\]
因此只要控制 $g$ 的上方，就能控制 \eqref{eq:win-positive-part} 的右边。若把左边擅自改成两个绝对值之和，就失去了这项性质。

对任意 $a>0$，指数函数的幂级数还给出
\begin{equation}
 \frac{u^ke^{-au}}{k!}\le a^{-k}\qquad(u\ge0),
 \label{eq:win-kernel-max}
\end{equation}
因为 $e^{au}\ge(au)^k/k!$。在好窗口内，由
$|g|+g\le2e^{\theta u}$ 和 \eqref{eq:win-kernel-max}，得到
\[
 \int_{W_U}(|g|+g)w_k
 \le \frac{2}{(\sigma-\theta)^k}\,\frac U2
 =\frac{U}{(\sigma-\theta)^k}.
\]
窗口外则用全局粗界 \eqref{eq:win-growth}。于是对每个好尺度 $U$、每个 $k\ge0$ 及每个 $t\in\mathbb R$，
\begin{equation}
 |A_k(\sigma+it)|+\operatorname{Re}A_k(\sigma)
 \le \frac{U}{(\sigma-\theta)^k}+2C_0T_{U,k},
 \label{eq:win-before-tail}
\end{equation}
其中
\[
 T_{U,k}=
 \int_{[0,\infty)\setminus W_U}
       \frac{u^ke^{-334u}}{k!}\,du.
\]
接下来的工作是控制这一尾积分。我们不能因为窗口外没有算术上界，就把那部分积分删去。

\section{用明确参数把积分核集中到窗口内}

对每个 $U\ge1$，取
\begin{equation}
 k=\lceil420U\rceil.
 \label{eq:win-k-choice}
\end{equation}
于是
\[
 420U\le k\le420U+1,\qquad k\ge U.
\]
导数阶数与窗口位置一起增大。相应权核的主要贡献落在窗口内部；下面的计算对窗口外剩余的贡献给出统一界。

先记录一个基本积分公式：
\begin{equation}
 \frac1{k!}\int_0^\infty u^ke^{-au}\,du
       =a^{-(k+1)}\qquad(a>0).
 \label{eq:win-gamma-integral}
\end{equation}
它从 $k=0$ 的指数积分出发，逐次分部积分即可得到；指数衰减使无穷远处的边界项为零。

在窗口右侧 $u\ge3U/2$，
\[
 e^{-334u}=e^{-320u}e^{-14u}
       \le e^{-21U}e^{-320u}.
\]
在窗口左侧 $0\le u\le U$，
\[
 e^{-334u}=e^{-348u}e^{14u}
       \le e^{14U}e^{-348u}.
\]
将两部分非负积分分别扩大到整个半轴，并使用
\eqref{eq:win-gamma-integral}，得到
\begin{equation}
 T_{U,k}\le
 \frac{e^{-21U}}{320^{k+1}}
 +\frac{e^{14U}}{348^{k+1}}.
 \label{eq:win-tail-two}
\end{equation}

为了与导数上界比较，把两项分别乘以 $\sigma^k=336^k$。由
$\log(1+x)\le x$（$x>-1$）和 \eqref{eq:win-k-choice}，
\begin{align*}
 336^k\frac{e^{-21U}}{320^{k+1}}
 &=\frac1{320}
   \exp\left(-21U+k\log\frac{21}{20}\right)\\
 &\le\frac1{320}\exp(-21U+k/20)
 \le\frac{e^{1/20}}{320}
 \le\frac e{320}.
\end{align*}
同样，
\begin{align*}
 336^k\frac{e^{14U}}{348^{k+1}}
 &=\frac1{348}
   \exp\left(14U+k\log\frac{28}{29}\right)\\
 &\le\frac1{348}\exp(14U-k/29)\\
 &\le\frac1{348}\exp(-14U/29)
 \le\frac1{348}.
\end{align*}
因此可以固定
\begin{equation}
 C_1=\frac e{320}+\frac1{348},
 \qquad
 T_{U,k}\le\frac{C_1}{\sigma^k}.
 \label{eq:win-tail-final}
\end{equation}
这里的 $C_1$ 与窗口位置 $U$ 无关。

令 $C_2=2C_0C_1$。由于 $0<\theta<1<\sigma$，
有 $\sigma^{-k}\le(\sigma-\theta)^{-k}$。
将 \eqref{eq:win-tail-final} 代入
\eqref{eq:win-before-tail}，得到本章的统一上界：
\begin{equation}
 \boxed{\;
 |A_k(\sigma+it)|+\operatorname{Re}A_k(\sigma)
 \le\frac{U+C_2}{(\sigma-\theta)^k},
 \qquad k=\lceil420U\rceil.
 \;}
 \label{eq:win-master-upper}
\end{equation}
它对每个好尺度 $U\ge1$ 和所有实数 $t$ 同时成立。

\section{在极点附近选出能够比较的两个圆盘}

现在反设存在 $\rho_0\in Z$，满足
$\operatorname{Re}\rho_0>\theta$。我们将选出一个实中心圆盘，其中没有极点；再竖直移动圆心，得到一个只含一个极点的圆盘。两个圆心都具有相同实部 $\sigma$，因此可以一起代入
\eqref{eq:win-master-upper}。

首先，由 $\Phi$ 在 $\operatorname{Re}s>2$ 全纯可知
\begin{equation}
 \operatorname{Re}\rho\le2<\sigma\qquad(\rho\in Z).
 \label{eq:win-poles-left}
\end{equation}
亚纯函数的极点在 $\mathcal H$ 内没有聚点，所以 $\mathcal H$ 中每个紧集只含有限个极点。

记 $R=\sigma-\theta$。闭圆盘
$\overline B(\sigma,R)$ 的实部下界为 $\theta>0$，故它紧含于 $\mathcal H$，其中的极点构成有限集。

如果该闭圆盘没有极点，取 $D=R$、$\rho_1=\rho_0$。此时 $B(\sigma,D)$ 无极点，且
\[
 0<\sigma-\operatorname{Re}\rho_1<D.
\]
如果闭圆盘中有极点，就在这有限多个极点中取一个离 $\sigma$ 最近的 $\rho_1$，并令
\[
 D=|\sigma-\rho_1|\le R.
\]
因为 $\sigma$ 不是极点，所以 $D>0$；最小性的定义保证开圆盘 $B(\sigma,D)$ 无极点。又因为正实轴没有极点，$\rho_1$ 非实，因而
\[
 D^2=(\sigma-\operatorname{Re}\rho_1)^2
             +(\operatorname{Im}\rho_1)^2
       >(\sigma-\operatorname{Re}\rho_1)^2.
\]
结合 \eqref{eq:win-poles-left}，仍得
$0<\sigma-\operatorname{Re}\rho_1<D$。

两种情形都给出了
\begin{equation}
 \begin{gathered}
 0<D\le\sigma-\theta,\qquad
 B(\sigma,D)\cap Z=\varnothing,\\
 \rho_1\in Z,\qquad
 0<a:=\sigma-\operatorname{Re}\rho_1<D.
 \end{gathered}
 \label{eq:win-free-disc}
\end{equation}
这一构造只需要检查紧圆盘中的有限个极点。

\subsection{竖移圆心并打破等距}

令
\[
 z_0=\sigma+i\operatorname{Im}\rho_1,
 \qquad |z_0-\rho_1|=a<D.
\]
选择
\begin{equation}
 0<h<\frac12\min\{D-a,\ \sigma-D\}.
 \label{eq:win-shift-margin}
\end{equation}
这是可能的，因为 $D<\sigma$。闭盘
$\overline B(z_0,D+h)$ 仍在 $\mathcal H$ 内，所以
\[
 P=Z\cap\overline B(z_0,D+h)
\]
是有限集。

当 $|t-\operatorname{Im}\rho_1|<h$，令 $z=\sigma+it$，则
\[
 |z-\rho_1|<h+a<D.
\]
此外，$B(z,D)$ 中的每个极点都属于 $P$，因为
\[
 |w-z_0|\le |w-z|+|z-z_0|<D+h.
\]
这样，一段小范围内所有可能影响最近距离的极点，都包含在同一个有限集合 $P$ 中。

对 $P$ 中两个不同点 $p,q$，两者到 $z=\sigma+it$ 的距离平方之差为
\begin{align*}
 &|\sigma+it-p|^2-|\sigma+it-q|^2\\
 &\quad=(\sigma-\operatorname{Re}p)^2
       -(\sigma-\operatorname{Re}q)^2
       +(\operatorname{Im}p)^2-(\operatorname{Im}q)^2
       -2t(\operatorname{Im}p-\operatorname{Im}q).
\end{align*}
如果两点虚部不同，至多一个 $t$ 会使此差为零。
如果两点虚部相同，则实部不同；由两点都严格位于 $\sigma$ 左侧可知
$(\sigma-\operatorname{Re}p)^2\ne
(\sigma-\operatorname{Re}q)^2$，所以没有这样的 $t$。

因此只需从开区间
$(\operatorname{Im}\rho_1-h,\operatorname{Im}\rho_1+h)$
中避开有限多个高度，就能使 $P$ 中各点到 $z$ 的距离互不相同。固定这样的 $t$ 与 $z$，取其中唯一的最近极点 $\rho$。它满足
\[
 0<d:=|z-\rho|<D.
\]
剩余各点的距离都严格大于 $d$；利用有限性，可以选到半径 $R_1$，使
\begin{equation}
 d<R_1<D,\qquad B(z,R_1)\cap Z=\{\rho\}.
 \label{eq:win-isolated-disc}
\end{equation}
若 $P$ 中只有 $\rho$，任取 $d<R_1<D$ 即可。又因
$R_1<D\le\sigma-\theta$，整个圆盘 $B(z,R_1)$ 都在 $\mathcal H$ 内。

到此已固定两个圆盘：$B(\sigma,D)$ 无极点，
$B(z,R_1)$ 只含一个单极点 $\rho$。所有选择均只依赖 $\Phi$、$\theta$ 及上述极点，与随后选取的好尺度 $U$ 无关。

\section{唯一最近极点给出导数下界}

选择一个共同的 Cauchy 半径
\begin{equation}
 d<r<R_1<D\le\sigma-\theta.
 \label{eq:win-cauchy-radius}
\end{equation}
在 $B(z,R_1)$ 内，减去 $\rho$ 处的主部，得到全纯函数
\[
 H(w)=\Phi(w)-\frac{c_\rho}{w-\rho}.
\]
这里在 $w=\rho$ 处取可去延拓；极点是单极点，保证只需减去这一项。

记
\[
 M_1=\max_{|w-z|=r}|H(w)|,\qquad
 M_0=\max_{|w-\sigma|=r}|\Phi(w)|.
\]
两者有限：第一个圆周位于 $H$ 的全纯圆盘内部，第二个圆周严格位于无极点圆盘 $B(\sigma,D)$ 内。Cauchy 导数估计给出，对每个 $k\ge0$，
\begin{align}
 \left|A_k(z)-\frac{c_\rho}{(z-\rho)^{k+1}}\right|
       &\le\frac{M_1}{r^k},\label{eq:win-pole-cauchy}\\
 |A_k(\sigma)|&\le\frac{M_0}{r^k}.
 \label{eq:win-free-cauchy}
\end{align}
第一式还使用了
\[
 \frac{(-1)^k}{k!}
 \left(\frac{c_\rho}{w-\rho}\right)^{(k)}_{\!w=z}
   =\frac{c_\rho}{(z-\rho)^{k+1}}.
\]

这两个估计把极点的贡献与其余全纯部分分开。极点项按 $d^{-k}$ 增长，而其余部分只需用较慢的 $r^{-k}$ 控制；严格距离差 $d<r$ 保证前者最终占优。

由 \eqref{eq:win-pole-cauchy} 和反三角不等式，
\[
 |A_k(z)|\ge\frac{|c_\rho|}{d^{k+1}}-\frac{M_1}{r^k}.
\]
另一方面，
\[
 \operatorname{Re}A_k(\sigma)
       \ge-|A_k(\sigma)|\ge-\frac{M_0}{r^k}.
\]
令
\begin{equation}
 b=\frac{|c_\rho|}{d}>0,\qquad B=M_0+M_1\ge0.
 \label{eq:win-lower-constants}
\end{equation}
于是得到与 \eqref{eq:win-master-upper} 左边完全相同的量的下界：
\begin{equation}
 |A_k(z)|+\operatorname{Re}A_k(\sigma)
       \ge\frac b{d^k}-\frac B{r^k}.
 \label{eq:win-master-lower}
\end{equation}
非零留数保证 $b$ 严格为正。这里无须要求复数留数本身是正实数。

\section{任意远处的好尺度产生矛盾}

现在才让好尺度 $U$ 向无穷远处移动。对每个这样的 $U\ge1$，取
$k=\lceil420U\rceil$。将
\eqref{eq:win-master-upper} 与
\eqref{eq:win-master-lower} 合并，有
\[
 \frac b{d^k}
 \le\frac B{r^k}+\frac{U+C_2}{(\sigma-\theta)^k}.
\]
两边乘以 $d^k$，并记
\[
 q=\frac dr\in(0,1),
\]
由于 $r\le\sigma-\theta$，可得
\begin{align}
 b
 &\le Bq^k+(U+C_2)
          \left(\frac d{\sigma-\theta}\right)^k\notag\\
 &\le(B+U+C_2)q^k
 \le(B+U+C_2)q^U.
 \label{eq:win-final-contradiction}
\end{align}
最后一步使用 $k\ge U$ 与 $0<q<1$。

所有常数 $b,B,C_2,q$ 都已固定，其中 $b>0$。
写 $q^U=e^{-\alpha U}$，$\alpha=-\log q>0$，即可看出
\[
 (B+U+C_2)q^U
       =(B+C_2+U)e^{-\alpha U}\longrightarrow0.
\]
指数衰减超过线性因子的增长。因此存在阈值 $U_*$，使全部 $U\ge U_*$ 的右边都小于 $b$。窗口假设保证仍能选到一个这样的好尺度，与
\eqref{eq:win-final-contradiction} 矛盾。

故不存在 $\rho_0\in Z$ 满足
$\operatorname{Re}\rho_0>\theta$。由于 $\mathcal H$ 中的 $\zeta$ 零点恰好成为 $\Phi$ 的极点，定理
\ref{thm:window-zeta} 得证。\hfill$\square$

这项证明只在选中的窗口内使用算术上界；窗口外的贡献由全局粗界和核尾估计负责。正实轴无极点的性质，则保证可以先找到无极点的实中心圆盘，再通过竖移圆心分离一个非实极点。两项条件各自承担了明确的步骤。

\section{两个反例说明条件不能随意删去}
取 $0<\theta<\beta<1$。若只假定一个实函数有单侧上界，而不检查实极点，那么
\[
g(u)=-e^{\beta u}\le e^{\theta u}
\]
甚至在全部 $u\ge0$ 都成立。但它的拉普拉斯变换为
\[
\Phi(s)=-\frac1{s-\beta},
\]
在 $\beta>\theta$ 有一个留数为负的实极点。因此，“单侧界必然排除所有极点”是错误的。本稿先证明 $\zeta$ 情形没有正实极点，才可以使用后面的圆盘选择。

另取 $g(u)=e^{\beta u}\cos u$，其拉普拉斯变换为
\[
\Phi(s)=\frac{s-\beta}{(s-\beta)^2+1},
\]
在 $\beta+i$ 和 $\beta-i$ 有极点。然而，对
$u=(2n+1)\pi$，有 $g(u)<0\le e^{\theta u}$，这样的好点能出现得任意远。甚至在每段余弦为负的固定长度区间上，上界都成立。
这仍不能排除极点，因为它没有提供本证明需要的、长度随 $U$ 增长的完整窗口 $[U,3U/2]$。

这两个例子说明，我们不能把证明中精确的窗口、实极点条件和导数组合省略后，仅凭“函数有上界”就宣布得到无零点结论。


% ===== finish.tex =====

\chapter{把排除边界推到二分之一}
\label{chap:finish}
\section{先排除临界线右侧}
现在将引用定理\ref{thm:arithmetic}与已经证明的窗口判据\ref{thm:window-zeta}接起来。

反设 $\rho\ne1$ 是一个零点，且
\[
\beta=\Rea\rho>\frac12.
\]
令
\[
a=\min\{\beta-\tfrac12,\tfrac12\}>0,
\qquad \theta=\frac12+\frac a2.
\]
因此 $0<\theta<1$ 且 $\theta<\beta$。算术定理的常数 $K>0$ 与 $\varepsilon_0>0$ 先固定，再取
\begin{equation}
\varepsilon=\min\left\{\frac{\varepsilon_0}{2},\frac{a}{4K}\right\}.
\label{eq:eps-final}
\end{equation}
这保证 $0<\varepsilon<\varepsilon_0$，并有
\[
\frac12+K\varepsilon\le\frac12+\frac a4
<\frac12+\frac a2=\theta<\beta.
\]
在算术定理提供的任意远窗口上，因 $x\ge1$，
\[
G(x)\le x^{1/2+K\varepsilon}\le x^\theta.
\]
窗口判据于是排除所有实部大于 $\theta$ 的零点，与 $\Rea\rho=\beta>\theta$ 矛盾。我们得到
\begin{equation}
\boxed{\zeta(s)\ne0\quad\text{当 }\Rea s>\tfrac12,\ s\ne1.}
\label{eq:right-free}
\end{equation}

\Expl{不必真的把 $\varepsilon$ 设为零。一个假想右侧零点与临界线之间有固定正距离；只要选择一次足够小但仍为正的 $\varepsilon$，就已经把它排除了。不同假想零点可以对应不同的小参数与窗口。}

\section{把函数方程写成可使用的形式}
定义完备化函数
\[
\mathcal Z(s)=\pi^{-s/2}\Gamma(s/2)\zeta(s).
\]
经典函数方程为亚纯恒等式
\begin{equation}
\mathcal Z(s)=\mathcal Z(1-s).
\label{eq:FE}
\end{equation}
例如，标准的 $\xi(s)=\tfrac12s(s-1)\mathcal Z(s)$ 满足
$\xi(s)=\xi(1-s)$；在 $s\ne0,1$ 处除去相等的多项式因子，再作亚纯延拓，就得到式\eqref{eq:FE}。标准公式见资料[4]。

Gamma 函数在 $0,-1,-2,\ldots$ 有单极点，在其他地方有限且非零。因而 $\Gamma(s/2)$ 的极点恰位于 $s=0,-2,-4,\ldots$。另外，$\pi^{-s/2}=\exp(-s\log\pi/2)$ 始终非零。Gamma 的这些性质见资料[5]。

\section{先核对 \texorpdfstring{$\zeta(0)$}{ζ(0)}，再处理左侧零点}
$\zeta$ 在 $0$ 全纯。由 $\Gamma(s/2)=2/s+O(1)$，
\[
\mathcal Z(s)=\frac{2\zeta(0)}s+O(1)\quad(s\to0).
\]
另一边，$\zeta(1-s)=-1/s+O(1)$，且
$\pi^{-(1-s)/2}\Gamma((1-s)/2)\to\pi^{-1/2}\Gamma(1/2)=1$。
函数方程给出 $2\zeta(0)=-1$，所以
\begin{equation}
\zeta(0)=-\frac12\ne0.
\label{eq:zeta-zero-value}
\end{equation}

现在设 $\rho$ 为非平凡零点且 $\Rea\rho<1/2$。由式\eqref{eq:zeta-zero-value}有 $\rho\ne0$；根据非平凡的限定，$\rho$ 也不是负偶数。因此 $\Gamma(\rho/2)$ 有限非零，$\mathcal Z$ 在 $\rho$ 有意义，并且
\[
\mathcal Z(\rho)=\pi^{-\rho/2}\Gamma(\rho/2)\zeta(\rho)=0.
\]
函数方程于是给出 $\mathcal Z(1-\rho)=0$。
但
\[
\Rea(1-\rho)=1-\Rea\rho>\frac12,
\qquad 1-\rho\ne1.
\]
由右侧非消失式\eqref{eq:right-free}，$\zeta(1-\rho)\ne0$。又因为
$\Rea((1-\rho)/2)>1/4$，Gamma 因子也有限非零，故
\[
\mathcal Z(1-\rho)=
\pi^{-(1-\rho)/2}\Gamma((1-\rho)/2)\zeta(1-\rho)\ne0,
\]
矛盾。左侧非平凡零点也被排除了。

\begin{center}
\begin{tikzpicture}[x=6cm,y=.65cm,>=Latex]
\fill[Teal!9] (.5,-2.5) rectangle (1.05,2.5);
\draw[->] (-.05,0)--(1.13,0) node[right] {$\Rea s$};
\draw[->] (0,-2.8)--(0,2.9) node[above] {$\Ima s$};
\draw[Teal,thick] (.5,-2.5)--(.5,2.5);
\node[below right] at (.5,0) {$1/2$};
\fill[red!65!black] (.25,2) circle (1.5pt);
\fill[red!65!black] (.75,-2) circle (1.5pt);
\draw[->,thick] (.25,2)--(.75,-2);
\node[above] at (.25,2) {$\rho=\beta+it$};
\node[below right] at (.75,-2) {$1-\rho=(1-\beta)-it$};
\node[align=center] at (.82,1.5) {右侧已被排除};
\end{tikzpicture}
\end{center}
图中是用于反证的假想位置。变换 $s\mapsto1-s$ 是绕 $(1/2,0)$ 的半周旋转，虚部也改变符号；这里没有把它画成同高度的横向反射。

\section{得到结论}
非平凡零点的实部既不能大于 $1/2$，也不能小于 $1/2$。所以
\[
\boxed{\Rea\rho=\frac12.}
\]
这完成了定理\ref{thm:RH}在引用定理\ref{thm:arithmetic}及明确列出的标准解析基础上的证明。



% ===== appendix.tex =====

\appendix
\chapter{深定理的来源和纸面证明的边界}
\label{chap:sources}
\section{可复查的源文件对应}
下表中的源码随制作包附上。完整路径保留在 \path{reference/source_index.json} 及相应源文件目录中；表内使用末级文件名，避免长路径遮住数学内容。

\begin{longtable}{@{}p{.31\textwidth}p{.64\textwidth}@{}}
\toprule
数学内容 & 原工程中的对应入口\\
\midrule
\endhead
平滑函数 $G$ 与族的定义 & \path{CharFamily.lean}；本稿只取模 $1$ 的主特征。\\
Mellin 变换及零点留数 & \path{poleData.lean} 及变换相关接口。正文亲自计算了主特征情形。\\
单侧窗口不等式 & \path{landauWindowIneq.lean}。关键组合保留了实轴导数项。\\
积分核的两侧尾部 & \path{landauTails.lean}。本稿取 $\delta=1/2$，把参数具体化为 $14,336,420$。\\
窗口排除零点 & \path{landauWindow.lean}、\path{ZeroFreeOfRpowWindow.lean}。本稿对 $\zeta$ 独立展开其解析论证，并删除不需要的实极点分支。\\
上同调长度窗口 & \path{ExistsPageN.lean}、\path{VwinOfRow.lean}、\path{TopMainZj.lean}。\\
宿主族及辅助条件 & \path{SupplyChi.lean}、\path{ZeroFreeAllOfTopG.lean}；更深构造由它们导入的文件证明。\\
对 $\varepsilon$ 统一的几何常数 & \path{DivinNOfHypV.lean} 中的 \path{FamilyClassConnected}，以及 \path{ShapeWin.lean}。\\
两类输入汇合 & \path{WindowInequalityOfArtinInputs.lean}、\path{LevelGmLeOnFamilyScales.lean}。\\
主项比较、误差与最终指数 & \path{ChTop.lean}、\path{ChCore.lean}、\path{SigALe.lean}、\path{CoreGm.lean}、\path{TopWin.lean}。\\
公开 RH 末端 & \path{RiemannHypothesisOfEveryCharacter.lean} 与公开主入口。本文用等价的经典完备化函数方程书写收尾，不要求读者先证明一般 Dirichlet GRH。\\
\bottomrule
\end{longtable}

\section{仍由原工程证明承担的工作}
下列部分在原工程中有对应证明与依赖，但没有在本稿中逐条转写：
\begin{enumerate}
\item theta 提升中的高度可积性、尾估计、局部函数、实位最高权计算及 Rallis 型非消失输入；
\item 自守表示到 Galois 表示的几何实现、同余极限与下降；
\item 基变换、局部相容性等输入导出上同调长度窗口的模性与估计论证；
\item 产生辅助素数及宿主族的可容许构造，并从同一族取得显式公式误差条件；
\item 产生相容几何链、秩条件和逐单位比较，最终给出式\eqref{eq:arith-main}所用的主项估计。
\end{enumerate}

所以，本稿中的算术输入有明确来源，但其来源本身仍需阅读全文才能成为完全自足的纸面证明。原作者的四页公告不是这些证明的替代文本；公开的形式核验认可，也不等于本文已经转写了所有依赖。

\section{本次实际完成的检查}
本次重新读取了前次保存的 TeX 和选定源码，并从固定镜像提交恢复完整原项目源码；恢复和单文件汇编详情见附录\ref{chap:source-delivery}。重构稿 TeX 的 SHA256 为
\begin{quote}\small
\path{072a7fe8afcb90f2758b689fb25b163cbf44b68217b164d93dada092c588ad23}
\end{quote}
恢复的原复核资料包 SHA256 为
\begin{quote}\small
\path{15ed5fd3ee92c73fa8beafe71f57da4d5aed01535c47ce3fdbee06fc941fa4bc}
\end{quote}
这两项与前次记录一致。本次还重新核查了当前社区声明，并对 $\zeta$ 专门化、量词顺序、常数吸收、单侧不等式、核尾估计、最近极点选择和函数方程的例外点进行了独立数学审阅。

本次未重新编译整个原 Lean 工程，也未把新稿全部再次形式化。先前小型 \path{EndpointReduction.lean} 的编译只验证带显式上游参数的末端推导；它与外部对完整公开工程的核验范围不同。

\section{与前一版视频的关系}
视频片尾展示的是 2026 年 11 月 21 日 20:55（北京时间）的 MIT 状态快照。新联合声明发布于同日 21:55（北京时间），报告了更进一步的共同体核验结果。时间较早的片尾记录应当按其标注时间理解。

本稿只讨论黎曼 $\zeta$ 函数。原工程还给出 Dirichlet $L$ 函数的相应结论；本稿没有把这个范围扩张为任意类型 $L$ 函数的广义黎曼猜想。

\chapter{从十四分卷到单文件源码与证明过程}
\label{chap:source-delivery}
\section{本轮使用的完整源码版本}
作者公告脚注列出 \path{lean.tar.gz.part01} 至 \path{lean.tar.gz.part14}。它们是同一个 gzip 归档的分卷，需要先按编号拼接，再解压；每个分卷并不是一个独立 Lean 模块。

本轮原站首分卷请求返回 HTTP 403，未重新取得十四份分卷字节。MIT 核验页公开提供源码镜像，并说明原源码与镜像相同，镜像后续另加比较器文件。本轮固定恢复社区仓库的首提交
\begin{quote}\small
\path{40e56bab026552e089d7afcfeff08d2defb9c005}
\end{quote}
的完整源码树，不把 GitHub 重新封装的 gzip 称为作者原始 gzip。先前保留的十四分卷 SHA256 清单同时附上，并标记本轮没有重新逐分卷校验。

实际恢复的项目包含 $175\,187$ 个 Lean 源文件，正文共 $1\,030\,531\,854$ 字节、$17\,374\,215$ 个物理行，另有三个项目配置文件。本轮将先前保留的 $88$ 个原源码及配置文件逐字节对照，全部一致。这一结果检查的是源码文本和来源，不能代替 Lean 内核对证明的检查。

\section{单文件汇编保留了什么}
交付文件 \path{RH_AllSources.lean} 按项目内部导入关系的拓扑顺序收录全部模块，在每个模块前后加入可识别的注释边界。各模块原始正文的字节保持不变，没有删除证明体，也没有用 \texttt{sorry} 或新增公理替代证明。

配套的 \path{RH_AllSources.index.jsonl} 记录模块名、原路径、正文的字节范围、行数和 SHA256；\path{RH_AllSources.dependencies.jsonl} 记录直接导入边。恢复工具读取汇编文件自身包含的模块边界、长度和哈希，抽取原文并放回原项目路径；恢复时不要求另有外部索引。因此保留了逐模块查阅和恢复编译工程的能力。直接导入图表示源文件依赖；它不等同于内核导出的最小定理依赖图。

\textbf{编译性质必须分清。}这个文件是无损源码汇编，尚未被改造为可独立编译的单一 Lean 模块。Lean 规定导入位于文件头；模块名还参与私有声明和作用域处理。把许多文件的正文拼接起来，会在正文中再次出现导入，并改变原来彼此隔离的局部环境。仅仅删除导入、给所有文件套一个命名空间，也不能保证保留原证明的含义。为此，本轮保留原文并提供可逆恢复，没有把文本合并的成功写成新单体证明的编译成功。

原项目的 Mathlib 与 Lean 工具链仍按 \path{lean-toolchain}、\path{lakefile.toml}、\path{lake-manifest.json} 固定。合并十四分卷所含源码，不会自动把 Lean 编译器、Mathlib 及全部外部依赖也变成同一个文件。

\clearpage
\section{三种证明文件各自说明什么}
\begin{longtable}{@{}p{.24\textwidth}p{.69\textwidth}@{}}
\toprule
文件种类 & 可以回答的问题\\
\midrule
\endhead
完整源码汇编 & 原作者写了哪些定义、定理和证明？某个模块的原文在哪里？能否无损恢复？\\
数学说明与关键源码 & 从目标声明到算术窗口、解析排零点和函数方程，各步骤如何相接？前提是什么？对应哪段代码？\\
内核证明项导出 & 展开后的声明与证明项是什么？独立检查器能否逐项接受？\\
\bottomrule
\end{longtable}

本轮实际导出的《证明过程\_数学说明与关键源码.txt》和《证明过程\_结构化索引.json》按十三个阶段组织说明，保留三十三个关键模块的原文、逐行编号及哈希。这个过程文件是源码级说明；没有把静态文本抽取冒充 Lean 执行时的逐 tactic 状态记录。

MIT 页面另公开提供完整内核导出 \path{solution.ndjson.zst}：压缩文件为 $7\,699\,933\,588$ 字节，解压为 $61\,570\,595\,788$ 字节。它与约一 GB 的 Lean 源码是不同层次的对象。本轮没有下载并复验该大文件；附带的来源记录给出其正式链接、哈希和公开核验入口，以便在具备相应资源的环境继续。不能把本文的过程说明当作该内核导出的替代品。

\section{如何沿证明过程继续阅读}
先从过程文件的 P01 核对最终命题，再沿 P03--P08 阅读窗口的算术来源，随后由 P09--P12 检查完整解析推导，最后在 P13 核对函数方程与特殊点。若只想理解正文的数学结构，可先读第\ref{chap:overview}章的记号表及第\ref{chap:window}章的上下界矛盾，再回到第\ref{chap:arithmetic}章逐项核对输入。

本次比上一版增加了上同调长度的有限部分和解释、宿主请求的量词、平方窗口缩到 $3/2$ 次方窗口、$401\to400$ 的截断余量、有限特征正交与 $4o$ 计数，以及统一阈值的拼接。尚需继续转写的是这些接口更深层的自守、Galois、几何和模性构造；源码完整收录不意味着这些构造已经全部成为人类文字证明。

\backmatter
\begin{thebibliography}{9}
\addcontentsline{toc}{chapter}{参考资料}
\bibitem{community}
Lean and Mathlib communities.
\emph{Lean verification of the Riemann Hypothesis}.
2026 年 11 月 21 日发布；本稿于北京时间 2026 年 11 月 22 日查阅。
\url{https://leanprover-community.github.io/blog/posts/riemann-hypothesis/}.

\bibitem{mit}
Andrew Sutherland.
\emph{Verification files for the Lean proof of the Riemann Hypothesis}.
本次打开的页面状态时间为 2026-11-21 12:55 UTC。
\url{https://math.mit.edu/~drew/grh/}.

\bibitem{announce}
Levent Alpöge.
四页 RH 公告及脚注中的 Lean 源归档。
\url{https://alpo.ge/remark.pdf}.
原始合并归档 SHA256：
\path{9631b4f020199e580ba86178cb56c6e9b09f690c2369a80dc4334aa93f49c7fc}.

\bibitem{functional}
NIST Digital Library of Mathematical Functions.
§25.4, Reflection Formulas，尤其式 25.4.3--25.4.4。
\url{https://dlmf.nist.gov/25.4}.

\bibitem{gamma}
NIST Digital Library of Mathematical Functions.
§5.2, Definitions，Gamma 函数的零点、极点与 Euler 常数。
\url{https://dlmf.nist.gov/5.2}.

\bibitem{leanmodules}
Lean Language Reference. Source Files and Modules.
\url{https://lean-lang.org/doc/reference/latest/Source-Files-and-Modules/}.

\bibitem{previous}
《黎曼猜想：Lean 证明的中文数学重构》，2026 年 11 月 21 日版本。
本次恢复其原始 TeX、复核资料包及选定源文件；本稿不是作者原文。
\end{thebibliography}
\end{document}
